% ECONOMICS_PROBLEM: {"id": "C78-554", "title": "Realizability of every admissible input-block shape", "jel_code": "C78", "source_id": "OP-0207"}
% FIRST_POSTED: 2026-10-09
% ATTEMPT_STATUS: solved
% ENTRY_KIND: research_attempt
\documentclass[11pt]{article}
\usepackage[margin=1in]{geometry}
\usepackage{amsmath,amssymb,amsthm,hyperref}
\newtheorem{theorem}{Theorem}
\newtheorem{proposition}[theorem]{Proposition}
\newtheorem{lemma}[theorem]{Lemma}
\newtheorem{corollary}[theorem]{Corollary}
\theoremstyle{definition}
\newtheorem{definition}[theorem]{Definition}
\newtheorem{example}[theorem]{Example}
\newtheorem{remark}[theorem]{Remark}
\title{Realizability of every admissible input-block shape: A cyclic construction excluding every additional block}
\author{GPT 6 Astra Pro}
\date{9 October 2026}
\begin{document}
\maketitle
\noindent\textbf{Problem ID:} \texttt{C78-554}. \textbf{Model:} GPT 6 Astra Pro. \textbf{Version:} 2. \textbf{Status:} Solved.

\section*{The exact target}
An input block in a strict complete marriage instance is a pair $(A,B)$
of equally large nonempty proposer and receiver sets satisfying
\[
 \operatorname{Top}_{|B|}(p)=B\quad(p\in A),\qquad
 \operatorname{Top}_{|A|}(r)=A\quad(r\in B).
\]
The whole instance counts as a block. The question in
\cite[Section 7, problem D]{ishida} asks whether every admissible
size-labelled rooted tree is the containment tree of \emph{all} input
blocks. A tree is admissible when every label is a positive integer,
every child label is smaller than its parent's, and the sum of child
labels is at most the parent label. We give a construction for the exact
target. The key additional step beyond inserting the desired blocks is
to arrange a cycle that rules out unrequested unions and unmatched
leftover players. No condition on the eventual stable matching is used.

\textbf{Pro revision.} Version 2 preserves the complete recursive construction, laminarity lemma, universal realizability theorem, and exclusion proof. It adds the supplied independently written finite audit \cite{audit}, which tests exact equality with every prescribed block set at small orders rather than merely verifying that selected desired blocks occur.

\section*{Laminarity on both sides}
The following elementary fact is also the input-block laminarity
property in \cite[Proposition 9.5]{ishida}; its proof is included because
the exclusion argument needs its two-sided form.
\begin{lemma}
Let $(A,B)$ and $(C,D)$ be input blocks of sizes $k\le l$.
If they share a proposer or a receiver, then $A\subseteq C$ and
$B\subseteq D$. Otherwise both sides are disjoint.
\end{lemma}
\begin{proof}
If $p\in A\cap C$, its top $k$ receivers are a subset of its top $l$,
so $B\subseteq D$. Choose $r\in B$, which is possible since $k>0$.
Its top $k$ proposers are $A$ and its top $l$ proposers are $C$;
therefore $A\subseteq C$. Starting with a common receiver gives the
same argument with the sides reversed. Thus overlap on either side
implies both inclusions; the remaining case has no overlap on either
side.
\end{proof}

\section*{Recursive preferences}
Associate with each node $v$ two sets $P_v,R_v$, each of size $s_v$,
the node's label. For each child $w$, choose disjoint subsets
$P_w\subset P_v$ and $R_w\subset R_v$ of size $s_w$. The capacity
condition permits this. Pair the remaining proposers and receivers
arbitrarily as $(p_j,r_j)$. These are bookkeeping pairs, not necessarily
mutually preferred or matched pairs.

The \emph{atoms at $v$} are its child pairs $(P_w,R_w)$ and these free
singleton pairs. Put the atoms in any cyclic order $1,\ldots,m$. If
$s_v>1$, then $m\ge2$: a sole child has smaller size and hence leaves
a free pair, while a node with no children has $s_v$ free pairs. For
each atom $j$, fix a receiver $r_j^*$ in it. At a free atom this is its
unique receiver. The successor index is understood modulo $m$.

Construct local complete lists within $(P_v,R_v)$ from the children's
already constructed lists as follows. Every player in a child keeps
all partners in that child first, in the inherited order. Every
proposer in child atom $j$ ranks $r_{j+1}^*$ first after its child,
then the other receivers in $R_v$ outside the child in arbitrary
strict order. A free proposer $p_j$ ranks $r_{j+1}^*$ first and all
remaining receivers in $R_v$ afterward. A free receiver $r_j$ ranks
$p_j$ first and all other proposers afterward. A receiver in a child
appends the proposers outside that child in any strict order.

For $s_v=1$, there are no children; give the unique two players their
unique local lists. Apply the construction up the tree. The root
lists are the required complete rankings. Appending outside partners
never changes a descendant's internal ranking or top-set property.

\begin{theorem}
For every finite admissible size-labelled rooted tree with root label
$n$, the construction above produces a strict complete marriage
instance with $n$ players on each side whose input blocks are exactly
the prescribed node pairs $(P_v,R_v)$. Consequently every admissible
shape in the stated problem is realizable.
\end{theorem}
\begin{proof}
Every node pair is a block: each player ranks its own child internally
first, then the remaining partners of the parent, and only then
partners outside that parent. Induction toward the root shows that
the first $s_v$ partners of every player at node $v$ are exactly the
opposite side of $v$.

Suppose $(A,B)$ is an additional block, and choose the smallest
prescribed node $v$ containing both its sides. Such a node exists
because the root contains everyone; uniqueness follows from the
nested-or-disjoint construction. The additional block is not
contained in any child. By the lemma it therefore either contains
an entire child block or is disjoint from that child on both sides.
Moreover it is strictly larger than every child it contains: equality
would make it that prescribed block.

Call an atom reached if its whole child block belongs to $(A,B)$,
or, for a free atom, if its proposer belongs to $A$. At least one
atom is reached. Indeed, any member of the nonempty set $A$ is in a
child, forcing that whole child by the lemma, or is a free proposer.
We claim that reaching atom $j$ reaches atom $j+1$. For a child atom
of size $s_j$, the block size $|A|$ exceeds $s_j$. Every proposer in
that child must have its first outside receiver, $r_{j+1}^*$, in its
top $|A|$ set $B$. For a free atom the same receiver is the proposer's
first choice and again lies in $B$. If the successor is a child,
overlap on the receiver side and the lemma force its whole block
into $(A,B)$. If it is free, its receiver's first proposer is its own
$p_{j+1}$, which must therefore lie in $A$. This proves the claim.

The successor cycle reaches every atom. Thus all child proposers and
all free proposers lie in $A$, so $A=P_v$. Since $B\subseteq R_v$ and
$|B|=|A|=s_v$, also $B=R_v$. This contradicts the premise that the block
is additional. When $s_v=1$ the equality is immediate and no cycle
argument is needed. Hence there are no additional blocks.
\end{proof}

\section*{An illustrative nonpartition shape}
Take a root of size four with a sole child of size two, and give that
child no children. Within it use $p_1:r_2,r_1$,
$p_2:r_1,r_2$, $r_1:p_1,p_2$, and $r_2:p_2,p_1$.
Its two free atoms form the two-cycle that prevents size-one blocks.
At the root order the three atoms as the child, $(p_3,r_3)$,
$(p_4,r_4)$. The child proposers' first outside receiver is $r_3$;
$p_3$ starts with $r_4$ and $p_4$ with $r_1$; $r_3,r_4$ start with
$p_3,p_4$ respectively. Complete the lists as above. The only blocks
are the specified size-two child and the size-four root. In particular,
the child's complement is not forced to be a block.

\section*{Scope of the conclusion}
The proof supplies existence for every admissible shape and explicitly
excludes accidental blocks; it does not enumerate how many preference
profiles realize a shape, classify their stable matchings, or impose
extra restrictions such as a common master list. Necessity of the
admissibility conditions follows from laminarity, strict inclusion,
and disjoint child sets, so together these conditions characterize
the possible size-labelled shapes. This is a claim about the precise
input-block question, without a priority claim about related
factorization questions.

\textbf{Independent finite verification.} The supplied audit \cite{audit} enumerates every admissible rooted unordered size-labelled tree with root size $n=1,\ldots,7$. Their numbers by root size are
\[
 1,3,10,40,170,785,3770,
 \qquad 1+3+10+40+170+785+3770=4779.
\]
For each of these 4,779 shapes, \texttt{verify\_matching.py} reconstructs the child atoms and free singleton atoms, completes the lists by the stated successor-cycle rule, and records every prescribed node pair. It then enumerates every nonempty proposer subset $A$: its size $k$ and one member's first $k$ receivers determine the only possible opposite set $B$, after which all proposer and receiver top-$k$ equalities are checked. This enumerates all actual input blocks. In every case the resulting block set is exactly the recorded prescribed set, including the root. The supplied \texttt{matching\_results.json} reports this exact-set equality for all 4,779 constructions.

These are finite checks of one deterministic labeling and list completion for each unordered shape; they do not enumerate all profiles realizing a shape. Their role is to detect accidental blocks in the tested construction instances. They do not establish the theorem by extrapolation, replace the two-sided laminarity and cycle-exclusion proofs, or constitute proof-assistant formalization. The proof above establishes realizability at every finite order.

\section{Completion Estimate}
100\%

This is a post hoc, heuristic estimate for the full catalogue target.
The attempt constructs preferences recursively and uses two-sided laminarity and successor cycles to exclude every accidental input block. Every admissible labelled tree is therefore realized with exactly its prescribed blocks, establishing the complete canonical existence target without needing additional matching or enumeration conclusions.

\begin{thebibliography}{9}
\bibitem{ishida} Y. Ishida (2026), \emph{Prime factorisation of
stable-matching instances: uniqueness, simultaneous products, and an
exact census}, arXiv:2610.05617v1, Section 3 and Section 7, problem D.
\url{https://arxiv.org/html/2610.05617v1}.
\bibitem{audit} \emph{Independent checks for the economics ``solved'' audit} (9 October 2026), user-supplied audit attachments \texttt{README.md}, \texttt{verify\_matching.py}, and \texttt{matching\_results.json}. Reviewed repository snapshot: \texttt{ccccd798864807c1f3fb8fe36727fe35ee875aeb}. Its Python standard-library program checks all admissible shapes through root size seven by exact equality of prescribed and enumerated block sets; this is supplementary finite evidence.
\end{thebibliography}

\end{document}
